% 演化映射的复合关系示意。曲线是同一条解在(t,x)平面中的图像，
% 三个时刻 r<s<t 上的状态依次记为 xi_r, xi_s, xi_t；不表示任何具体方程的数值。
\definecolor{OdeBlue}{HTML}{7188A3}\definecolor{OdeBlueFill}{HTML}{E5EAF0}
\definecolor{OdeGreen}{HTML}{829C8B}\definecolor{OdeYellow}{HTML}{B4A16F}
\definecolor{OdeRed}{HTML}{A7787D}\definecolor{OdeInk}{HTML}{3B4252}
\definecolor{OdeMuted}{HTML}{6B7280}\definecolor{OdeStroke}{HTML}{4C4D4F}
\begin{tikzpicture}[line width=.6pt,font=\fontsize{8}{10}\selectfont,>=Stealth,
  dot/.style={circle,fill=OdeInk,inner sep=1.3pt},
  note/.style={OdeMuted,font=\fontsize{7.5}{9}\selectfont}]
\draw[->,OdeStroke] (0,0) -- (98mm,0) node[below left,OdeInk] {$t$};
\draw[->,OdeStroke] (0,0) -- (0,40mm) node[below left,OdeInk] {$\bm x$};
\draw[OdeBlue,line width=1.1pt] plot[smooth,tension=.7] coordinates
  {(6mm,8mm) (18mm,13mm) (32mm,24mm) (48mm,27mm) (64mm,20mm) (78mm,25mm) (90mm,32mm)};
\node[dot] (a) at (18mm,13mm) {};
\node[dot] (b) at (48mm,27mm) {};
\node[dot] (c) at (78mm,25mm) {};
\foreach \p/\x in {a/18,b/48,c/78}{\draw[OdeMuted,dashed,line width=.4pt] (\p) -- (\x mm,0);}
\node[below,OdeInk] at (18mm,0) {$r$};
\node[below,OdeInk] at (48mm,0) {$s$};
\node[below,OdeInk] at (78mm,0) {$t$};
\node[below right,OdeInk] at (a) {$\bm\xi_r$};
\node[above,OdeInk,yshift=1pt] at (b) {$\bm\xi_s$};
\node[below right,OdeInk] at (c) {$\bm\xi_t$};
\draw[->,OdeGreen,line width=.8pt] (a) to[bend left=38] node[above left,pos=.45] {$\bm\varphi_{s,r}$} (b);
\draw[->,OdeGreen,line width=.8pt] (b) to[bend left=38] node[above right,pos=.55] {$\bm\varphi_{t,s}$} (c);
\draw[->,OdeRed,line width=.8pt] (a) to[bend right=16] node[below,pos=.5,yshift=-1pt] {$\bm\varphi_{t,r}=\bm\varphi_{t,s}\circ\bm\varphi_{s,r}$} (c);
\node[note,anchor=west] at (4mm,36mm) {同一条解经过三个时刻};
\end{tikzpicture}
